\documentclass[10pt,a4paper]{amsart}
\usepackage[margin=19mm]{geometry}
\usepackage{amsmath,amssymb,mathtools}
\usepackage[scr=rsfs,cal=euler]{mathalfa}
\usepackage{xfrac}
\usepackage[unicode,hidelinks,bookmarks=false]{hyperref}
\newcommand{\ie}{i.e.\ }
\newcommand{\cf}{cf.\ }
\newcommand{\resp}{respectively\ }
\newcommand{\Subsection}{\subsection}
\providecommand{\nobreakdash}{\nobreak\hbox{-}}
\setlength{\emergencystretch}{2em}
\multlinegap0pt
\usepackage{amssymb}
\usepackage{xcolor}
\newtheorem{theorem}{Theorem}
\def\h{\mathbb H}
\def\r{\mathbb R}
\title{Geometric aspects of the curve shortening flow in the hyperbolic plane}
\author{Ivan Krznari\'c, Rafael L\'opez}
\date{Source: arXiv:2605.13023v1 (CC BY 4.0)}
\begin{document}
\maketitle

\noindent\emph{Source and licence.} Geometric aspects of the curve shortening flow in the hyperbolic plane. Version arXiv:2605.13023v1 (CC BY 4.0), distributed by arXiv under the Creative Commons Attribution 4.0 International licence, http://creativecommons.org/licenses/by/4.0/, as recorded in the arXiv licence metadata for this exact version and shown on the abstract page https://arxiv.org/abs/2605.13023v1. Full licence notice: \texttt{licenses/arxiv-2605.13023-CC-BY-4.0.txt}.\par\smallskip

\noindent\textit{Editorial scope.} Counted unit: the article's theorem t42 (source0.tex lines 486--490). The article's complete original proof of it is delivered with it (source0.tex lines 492--545, one \texttt{proof} environment of 54 lines). No mathematics is written, repaired or re-derived: the statement and the proof are the article's own lines, in the article's own notation, and no retained line is substituted or re-flowed.  Retained uncounted as the source-local context the proof needs: lines 457--463.  Nothing was omitted from any retained span, so the package carries no omission record.  No retained line contains a citation, so the wrapper carries no bibliography and no bibliography entry.  No retained line contains a figure environment, an image-inclusion command or drawing code, so no image is delivered and the package declares no asset dependency.  Editorial formatting only, in addition: an AMS \texttt{amsart} preamble that restates the article's own declarations and macros used by the retained lines, namely the environments \texttt{theorem} on separate counters, together with the standard packages those lines need.  The retained environments print in this document as Theorem 1 in order of appearance, whereas the article prints the same items with the numbering of its own full document.  The complete native source of the article as distributed in the arXiv e-print for this exact version is bundled unchanged as \texttt{sources/200474/source0.tex}.
\begin{equation}\label{evo}
\left\{
\begin{split}
p_\tau &= p p_{\varphi \varphi} - \frac{1}{2}p_{\varphi}^2 + 2p^2 - 2p \\
 q_\tau&= p(q_{\varphi \varphi} + 4q) - 2 q
\end{split}\right.
\end{equation}

\begin{theorem}\label{t42}
Let $\{\gamma_t\colon t \in J\}$ be a convex solution to the CSF in $\h^2$ such that the corresponding pressure function $p$ can be written as 
$$p(t, \varphi) = a(t) + b(\varphi),$$
for some functions $a,b$. Then the solution   is a  family of shrinking circles, horocycles, equidistant lines or solitons.
\end{theorem}

\begin{proof}
Equations \eqref{evo} become
\begin{equation}\label{evo2}
\left\{
\begin{split}
a' &= (a+b)b'' - \frac{1}{2}b'^2 + 2(a+b)^2 - 2(a+b) \\
0&= (a+b)(b'''+4b')-2b',
\end{split}\right.
\end{equation}
where $'$ denotes, in each case, the corresponding derivative with respect to the variable $t$ or $\varphi$. 
 If the function $a(t)=a$ is constant, then at any time $t$, the curves $\gamma(t, \cdot)$ and $\gamma(0, \cdot)$ have the same curvature function, namely $\kappa(t, \cdot) = \sqrt{a + b(\cdot)} = \kappa(0, \cdot)$. 
Because two curves have the same curvature, there exists an isometry of the ambient space $\Phi(t)$ such that $\gamma(t, \cdot) = \Phi(t)\gamma(0, \cdot)$. Since $t$ is an arbitrary time, the solution evolves by isometries,  that is, $\gamma$ is a soliton.  
 
Now assume that $a$ is not a constant function. We distinguish two cases.
\begin{enumerate}
\item Assume   $b(\varphi) $ is a constant function. Without loss of generality, we can assume that this constant is $0$. Then $p(t, \varphi) = a(t)$.  The second equation in \eqref{evo2} is trivial and the first one   reduces to  
$$a' = 2a^2 - 2a.$$ 
We now solve this equation. This equation can be written as
$$\frac{da}{a(a-1)}=2\, dt,$$
obtaining
$$a(t)=\frac{1}{1+C e^{2t}},\quad C\in\r.$$
Since the function $a$ is positive, we have three types of equations depending on the sign of the integration constant $C$:
\begin{equation}\label{a3}
a(t) = \left\{\begin{array}{lll} 
 &\frac{1}{1-Ce^{2t}},  & t<-\frac12\log(C), C>0\\ 
  & 1 &t\in\r\\
  &\frac{1}{1+Ce^{2t}},&t\in\r,  C>0.
  \end{array}\right.
\end{equation}
\begin{enumerate}
\item Case $a(t) =  \frac{1}{1-Ce^{2t}}$, where $t\in (-\infty,-\frac12\log(C))$ and $C>0$. Let $c=\frac12\log(C)$. Then $C=e^{2c}$. Define
$$r(t)=\cosh^{-1}(e^{-(t+c)}).$$
 Then we have 
$$\kappa(t, \varphi) = \sqrt{p(t, \varphi)} = \sqrt{a(t)} = \frac{1}{\tanh(r(t))}.$$ 
This implies that the  corresponding solution is a  family of shrinking circles.  
 \item Case  $a(t)=1$, $t\in\r$. Then $\kappa=1$. Thus $\gamma(t,\cdot)$ is a horocycle for all $t$. 

\item Case $a(t) =  \frac{1}{1+Ce^{2t}}$, where $t\in \r$ and $C>0$. Again, if $c=\frac12\log(C)$, let   
$$r(t)=\sinh^{-1}(e^{t+c}).$$
 Then
$$\kappa(t, \varphi) = \sqrt{p(t, \varphi)} = \sqrt{a(t)} = \frac{1}{\cosh(r(t))}.$$ 
In consequence, the solution is a family of equidistant lines. Notice that $k<1$.
 \end{enumerate}
 
 
\item Assume $b$ is not a constant function.   The second equation of \eqref{evo2} is  
\begin{equation}\label{e3}
(a+b)(b''' + 4b') - 2b' = 0.
\end{equation}
If $b'''+4b'=0$ identically, then \eqref{e3} implies $b'=0$. Thus $b$ is a constant function, which it is a contradiction. Thus  $b'''+4b'\not=0$.  From \eqref{e3}, we have 
 $$a(t) = \frac{2b'(\varphi)}{b''' (\varphi)+ 4b'(\varphi)} - b(\varphi).$$ 
The left-hand side of the equation above depends only on $t$ while the right-hand side depends only on $\varphi$. In particular, the function $a$ is constant. This contradiction proves that this case is not possible.  
\end{enumerate}
\end{proof}

\end{document}
